% =====================================================================
%  BÖLÜM 7: Büyük Sayılar Kanunları
%  Kaynak: Feller Vol.I-II; Billingsley; Durrett; Klenke
% =====================================================================
\chapter{Büyük Sayılar Kanunları}
\label{chp:buyuk_sayilar}

\bolumalinti{%
  Büyük sayılar kanunu olasılığın en derin gerçeğidir:\\
  belirsizlik çoğunlukla sönümlenir.%
}{Richard Durrett, \textit{Probability: Theory and Examples}}

% ---- Kazanımlar (TYMM) ----
\begin{kazanimlar}
Bu bölümün sonunda öğrenci:
\begin{itemize}
  \item Neredeyse kesin, olasılıkla, dağılımla ve $L^p$ yakınsamalarını tanımlayabilir;
        aralarındaki hiyerarşik ilişkileri ispatlayabilir.
  \item Borel-Cantelli lemmalarını (her ikisini) ispatlayabilir ve uygulayabilir.
  \item Zayıf Büyük Sayılar Kanunu'nu (ZBSK) Chebyshev yöntemiyle ispatlayabilir;
        Helly seçim teoremi ile genel versiyonunu açıklayabilir.
  \item Güçlü Büyük Sayılar Kanunu'nu (GBSK) Kolmogorov yaklaşımıyla ispatlayabilir.
  \item Kolmogorov'un 0-1 Kanunu'nu ifade edip ispatlayabilir;
        "erillik" ve olasılık-0 veya 1 sonuçlarını tanıyabilir.
  \item Büyük sapma üst sınırlarını (Chernoff sınırı) türetebilir ve uygulayabilir.
\end{itemize}
\end{kazanimlar}

% ---- Motivasyon ----
\begin{motivasyon}
Bir para atışında yazı gelme oranı $n$ büyüdükçe $1/2$'ye yaklaşır — bu sezgi,
matematiksel olarak tam olarak ne anlama gelir? "Yaklaşmak" birden fazla biçimde
tanımlanabilir: \emph{her} örnek yolu için mi, yoksa çoğunlukla mı?

Büyük Sayılar Kanunları bu soruyu yanıtlar. Zayıf versiyon (ZBSK) olasılıkla
yakınsamayı, güçlü versiyon (GBSK) ise neredeyse kesin yakınsamayı garantiler.
Aradaki fark ince ama derin: GBSK, "uzun vadede ortalama gerçek beklentiye eşit
olur" ifadesinin en kesin matematiksel anlatımıdır.

\medskip
\textbf{Köprü.} Bölüm~\ref{chp:beklenti}'deki Chebyshev eşitsizliği ZBSK'nın
kalbi; Bölüm~\ref{chp:olcum_teorisi}'deki Monoton Yakınsama Teoremi GBSK'nın
temelidir. Bölüm~\ref{chp:merkezi_limit}'te yakınsama \emph{hızı} araştırılacak.
\end{motivasyon}

% =====================================================================
\section{Yakınsama Türleri}
\label{sec:yaklinsama_turleri}
% =====================================================================

\begin{motivasyon}[Neden Birden Fazla Yakınsama Türü?]
$X_n \to X$ demek nedir? Dizinin her terimi bir rasgele değişken.
Analiz derslerinin "nokta yakınsaması" (her $\omega$ için) çok kuvvetli;
"düzgün yakınsama" daha da kuvvetli. Olasılık teorisi bu ikisinin yanı sıra
daha zayıf ama yeterince güçlü iki kavram sunar: olasılıkla ve dağılımla yakınsama.
\end{motivasyon}

\begin{tanim}[Neredeyse Kesin Yakınsama]
\label{def:nk_yaklinsama}
$\{X_n\}$ rasgele değişkenler dizisi. $X_n \xrightarrow{\text{n.k.}} X$
(\textbf{neredeyse kesin yakınsama}, almost sure convergence) denir, eğer
\[
  P\!\left(\left\{\omega : \lim_{n\to\infty} X_n(\omega) = X(\omega)\right\}\right) = 1.
\]
Eşdeğer: $P(|X_n - X| > \varepsilon \text{ sonsuz çok } n \text{ için}) = 0$ her $\varepsilon > 0$'da.
\end{tanim}

\begin{tanim}[Olasılıkla Yakınsama]
\label{def:p_yaklinsama}
$X_n \xrightarrow{P} X$ (\textbf{olasılıkla yakınsama}, convergence in probability) denir, eğer
her $\varepsilon > 0$ için
\[
  \lim_{n\to\infty} P(|X_n - X| > \varepsilon) = 0.
\]
\end{tanim}

\begin{tanim}[Dağılımla Yakınsama]
\label{def:d_yaklinsama}
$X_n \xrightarrow{d} X$ (\textbf{dağılımla yakınsama}, convergence in distribution) denir, eğer
$F_X$ sürekli olan her $x$'te
\[
  \lim_{n\to\infty} F_{X_n}(x) = F_X(x).
\]
\end{tanim}

\begin{tanim}[$L^p$ Yakınsaması]
\label{def:lp_yaklinsama}
$p \geq 1$ ve $E[|X_n|^p], E[|X|^p] < \infty$ ise $X_n \xrightarrow{L^p} X$ denir, eğer
\[
  \lim_{n\to\infty} E[|X_n - X|^p] = 0.
\]
\end{tanim}

% =====================================================================
\section{Yakınsama Türleri Arasındaki İlişkiler}
\label{sec:yaklinsama_iliskileri}
% =====================================================================

\begin{teorem}[Yakınsama Hiyerarşisi]
\label{thm:yaklinsama_hiyerarsi}
Aşağıdaki çıkarımlar geçerlidir:
\[
  X_n \xrightarrow{\text{n.k.}} X \;\Rightarrow\; X_n \xrightarrow{P} X \;\Rightarrow\; X_n \xrightarrow{d} X,
\]
\[
  X_n \xrightarrow{L^p} X \;\Rightarrow\; X_n \xrightarrow{P} X.
\]
Hiçbir ok ters yönde genel olarak geçerli değildir.
\end{teorem}

\begin{proof}
\textbf{n.k.\ $\Rightarrow$ P:} Sabit $\varepsilon > 0$ için
$\limsup_n \{|X_n - X| > \varepsilon\} \subseteq \{\omega : X_n(\omega) \not\to X(\omega)\}$.
n.k.\ yakınsama: sağ tarafın olasılığı sıfır. Ölçünün monotonluğu:
$P(|X_n-X| > \varepsilon) \leq P(\bigcup_{k\geq n}\{|X_k-X|>\varepsilon\})$,
bu da $n\to\infty$'da sıfıra gider (azalan olaylar dizisi, Bölüm~\ref{chp:olcum_teorisi}).

\textbf{P $\Rightarrow$ d:} $\varepsilon > 0$ ve $x$ süreklilik noktası alın.
$F_{X_n}(x) = P(X_n \leq x) \leq P(X \leq x+\varepsilon) + P(|X_n - X| > \varepsilon)$.
$n\to\infty$: $\limsup F_{X_n}(x) \leq F_X(x+\varepsilon)$.
Benzer şekilde alt sınır kurulur; $\varepsilon\to 0$ ile $F_{X_n}(x) \to F_X(x)$.

\textbf{$L^p$ $\Rightarrow$ P:} Markov ($|X_n-X|^p \geq \varepsilon^p \mathbf{1}_{|X_n-X|>\varepsilon}$):
$P(|X_n-X|>\varepsilon) \leq E[|X_n-X|^p]/\varepsilon^p \to 0$.
\end{proof}

\begin{karsiornek}[Okların Terslenmediğine Örnekler]
\label{ex:yaklinsama_karsiornek}

\textbf{P $\not\Rightarrow$ n.k.:} $\Omega = [0,1]$, $P = \lambda$. "Gezgin blok":
$X_n = \mathbf{1}_{[(k-1)/m, k/m]}$ ($n = m(m-1)/2 + k$, $1 \leq k \leq m$).
$P(|X_n| > 0) = 1/m \to 0$ ama her $\omega$ için sonsuz çok $n$'de $X_n(\omega) = 1$.

\textbf{d $\not\Rightarrow$ P:} $X_n \equiv X \sim \Normal(0,1)$ ama $Y = -X \sim \Normal(0,1)$.
$X_n \xrightarrow{d} Y$ ama $X_n - Y = 2X$ sıfıra olasılıkla yakınsamaz.
\end{karsiornek}

\begin{teorem}[Slutsky Teoremi]
\label{thm:slutsky}
$X_n \xrightarrow{d} X$ ve $Y_n \xrightarrow{P} c$ (sabit) ise:
\begin{enumerate}[(i)]
  \item $X_n + Y_n \xrightarrow{d} X + c$,
  \item $X_n Y_n \xrightarrow{d} cX$,
  \item $X_n / Y_n \xrightarrow{d} X/c$ ($c \neq 0$).
\end{enumerate}
\end{teorem}

\begin{proof}
\textit{(i):} $P(X_n + Y_n \leq x) \leq P(X_n \leq x - c + \varepsilon) + P(|Y_n - c| > \varepsilon)$.
$n\to\infty$: $\limsup \leq F_X(x-c+\varepsilon)$; $\varepsilon\to 0$: $F_{X+c}(x)$.
Alt sınır benzer şekilde; eşitlik sağlanır.
\end{proof}

\begin{teorem}[Sürekli Haritalama Teoremi]
\label{thm:surekli_haritalama}
$g : \mathbb{R} \to \mathbb{R}$ sürekli ve $X_n \xrightarrow{d} X$ (veya $\xrightarrow{P}$ veya $\xrightarrow{\text{n.k.}}$)
ise $g(X_n) \xrightarrow{d} g(X)$ (veya $\xrightarrow{P}$, $\xrightarrow{\text{n.k.}}$ sırasıyla).
\end{teorem}

\begin{proof}
n.k.\ durumu: $X_n(\omega) \to X(\omega)$ ve $g$ sürekli $\Rightarrow$ $g(X_n(\omega)) \to g(X(\omega))$
($P = 1$ kümesi üzerinde).
Dağılımla yakınsama için Portmanteau teoremi kullanılır (ayrıntı: Billingsley~\cite{Billingsley1995}).
\end{proof}

\begin{mikroozet}
Kuvvetli sıralaması: n.k.\ $\Rightarrow$ P $\Rightarrow$ d; $L^p$ $\Rightarrow$ P.
Tersleri genel olarak yanlış. Slutsky: $X_n \xrightarrow{d} X$, $Y_n \xrightarrow{P} c \Rightarrow X_n+Y_n \xrightarrow{d} X+c$.
\end{mikroozet}

% =====================================================================
\section{Borel-Cantelli Lemmaları}
\label{sec:borel_cantelli}
% =====================================================================

\begin{kesif}
$\{A_n\}$ olaylar dizisi. "$A_n$ sonsuz çok gerçekleşir" olayını
Bölüm~\ref{chp:onbilgiler}'deki limsup dili ile nasıl ifade ederiz?
Eğer $\sum P(A_n) < \infty$ ise bu olayın gerçekleşme olasılığı hakkında
ne söyleyebiliriz?
\kesifcevap{``$A_n$ sonsuz çok gerçekleşir'' $= \limsup A_n = \bigcap_{n=1}^\infty \bigcup_{k=n}^\infty A_k$.
  Borel-Cantelli I: $\sum P(A_n)<\infty \Rightarrow P(\limsup A_n)=0$.
  Sezgi: toplam olasılık sonludur, dolayısıyla sonsuz kez gerçekleşme imkânsızdır (n.k.).}
\end{kesif}

\begin{tanim}[Sonsuz Sık Olay]
$\{A_n\}$ olaylar dizisi için $\limsup_n A_n = \bigcap_{n=1}^\infty \bigcup_{k=n}^\infty A_k$.
Bu olay "\textbf{$A_n$ sonsuz çok gerçekleşir}" (i.o.: infinitely often) diye okunur.
\end{tanim}

\begin{teorem}[Borel-Cantelli Birinci Lemması]
\label{thm:bc1}
$\{A_n\} \subseteq \mathcal{F}$ ve $\sum_{n=1}^\infty P(A_n) < \infty$ ise
$P(\limsup_n A_n) = 0$.
\end{teorem}

\begin{proof}
$B_n = \bigcup_{k=n}^\infty A_k$ için monoton azalma $B_n \downarrow \limsup_n A_n$.
\[
  P(\limsup_n A_n) = \lim_{n\to\infty} P(B_n) \leq \lim_{n\to\infty} \sum_{k=n}^\infty P(A_k) = 0,
\]
çünkü $\sum P(A_n) < \infty$ olduğundan kuyruk toplamı sıfıra gider.
\end{proof}

\begin{teorem}[Borel-Cantelli İkinci Lemması]
\label{thm:bc2}
$\{A_n\}$ \textbf{bağımsız} olaylar ve $\sum_{n=1}^\infty P(A_n) = \infty$ ise
$P(\limsup_n A_n) = 1$.
\end{teorem}

\begin{proof}
$P(\limsup_n A_n)^c = P(\liminf_n A_n^c)$.
$P\!\left(\bigcap_{k=n}^m A_k^c\right) = \prod_{k=n}^m (1-P(A_k))$ (bağımsızlık).
$1 - x \leq e^{-x}$: $\prod_{k=n}^m (1-P(A_k)) \leq \exp\!\left(-\sum_{k=n}^m P(A_k)\right) \to 0$
($\sum P(A_k) = \infty$ olduğundan). Yani $P(\bigcap_{k=n}^\infty A_k^c) = 0$ her $n$ için.
Sayılabilir birleşim: $P(\liminf A_n^c) = P(\bigcup_n \bigcap_{k\geq n} A_k^c) = 0$.
\end{proof}

\begin{hataanalizi}
İkinci lemma için bağımsızlık şartı çıkarılamaz.
Örnek: $A_n = A$ sabit olay, $P(A) = 1/2$, $\sum P(A_n) = \infty$.
Ama $P(\limsup A_n) = P(A) = 1/2 \neq 1$.
Bağımsızlık olmadan $\sum P(A_n) = \infty$ yalnızca üst sınır verir, kesin sonuç değil.
\end{hataanalizi}

\begin{ornek}[Para Atışı]
$A_n = \{n\text{'inci atışta arka}\ k\ \text{kez art arda}\}$.
$P(A_n) = 1/2^k$. $\sum P(A_n) = \infty$ ve $\{A_n\}$ bağımsız.
İkinci Borel-Cantelli: herhangi bir $k$ için art arda $k$ arka gelme n.k.\ sonsuz kez gerçekleşir.
\end{ornek}

% =====================================================================
\section{Zayıf Büyük Sayılar Kanunu}
\label{sec:zbsk}
% =====================================================================

\begin{tanim}[Örneklem Ortalaması]
$X_1, X_2, \ldots$ rasgele değişkenler; $\bar{X}_n = \frac{1}{n}\sum_{i=1}^n X_i$.
\end{tanim}

\begin{teorem}[Chebyshev ZBSK]
\label{thm:chebyshev_zbsk}
$X_1, X_2, \ldots$ çiftler arası korelasyonsuz (pairwise uncorrelated),
$E[X_i] = \mu$ ve $\mathrm{Var}(X_i) \leq M < \infty$ ise
$\bar{X}_n \xrightarrow{P} \mu$.
\end{teorem}

\begin{proof}
$E[\bar{X}_n] = \mu$. Korelasyonsuzluktan:
$\mathrm{Var}(\bar{X}_n) = \frac{1}{n^2}\sum_{i=1}^n \mathrm{Var}(X_i) \leq \frac{M}{n}$.
Chebyshev: $P(|\bar{X}_n - \mu| > \varepsilon) \leq \mathrm{Var}(\bar{X}_n)/\varepsilon^2 \leq M/(n\varepsilon^2) \to 0$.
\end{proof}

\begin{teorem}[Genel ZBSK (Kolmogorov)]
\label{thm:zbsk_genel}
$X_1, X_2, \ldots$ \textbf{i.i.d.}, $E[|X_1|] < \infty$, $E[X_1] = \mu$ ise
$\bar{X}_n \xrightarrow{P} \mu$.
\end{teorem}

\begin{proof}[Kırpma Yöntemi]
$\mu = 0$ (genel durum ötelemeli indirgenir). Sabit $M > 0$ için kırpılmış rasgele değişken:
$Y_i = X_i \mathbf{1}_{|X_i| \leq M}$.
$E[Y_i] = E[X_i \mathbf{1}_{|X_i|\leq M}] \to E[X_i] = 0$ ($M\to\infty$, DCT).
$\mathrm{Var}(Y_i) \leq E[Y_i^2] \leq ME[|Y_i|] \leq M \cdot E[|X_1|] < \infty$.

$\bar{X}_n - \bar{Y}_n = \frac{1}{n}\sum(X_i - Y_i) = \frac{1}{n}\sum X_i \mathbf{1}_{|X_i|>M}$.
Markov: $P(|\bar{X}_n - \bar{Y}_n| > \varepsilon) \leq \frac{1}{\varepsilon}E[|X_1|\mathbf{1}_{|X_1|>M}] \to 0$
($M\to\infty$, DCT).
$\bar{Y}_n \xrightarrow{P} E[Y_1] \approx 0$ (Chebyshev, $M$ büyük); ikisini birleştirince $\bar{X}_n \xrightarrow{P} 0$.
\end{proof}

\begin{ornek}[Monte Carlo Entegrasyonu]
$g : [0,1] \to \mathbb{R}$ integrallenebilir. $U_1, \ldots, U_n \sim \Uniform(0,1)$ bağımsız.
$\frac{1}{n}\sum g(U_i) \xrightarrow{P} E[g(U)] = \int_0^1 g(x)\,dx$.
ZBSK, Monte Carlo yönteminin matematiksel garantisidir.
\end{ornek}

\begin{kritiksorgu}
ZBSK şunu söyler: $P(|\bar{X}_n - \mu| > \varepsilon) \to 0$.
Bu, $\bar{X}_n$'nin hiçbir zaman $\mu$'dan uzak olmadığı anlamına mı gelir?
Hayır — yalnızca büyük sapmalar giderek daha düşük olasılıklı hale gelir.
$\bar{X}_n$'nin $\mu$'dan uzak olduğu anların kümesi boş olmayabilir;
ama bu anların "oranı" sıfıra yaklaşır.
\end{kritiksorgu}

% =====================================================================
\section{Güçlü Büyük Sayılar Kanunu}
\label{sec:gbsk}
% =====================================================================

\begin{teorem}[Kolmogorov GBSK]
\label{thm:gbsk}
$X_1, X_2, \ldots$ \textbf{i.i.d.}, $E[|X_1|] < \infty$, $E[X_1] = \mu$ ise
$\bar{X}_n \xrightarrow{\text{n.k.}} \mu$.
\end{teorem}

\begin{proof}
$\mu = 0$ varsayalım. Dört-moment argümanı (i.i.d., sonlu dördüncü moment varsayarak):

\textbf{Adım 1 (Dördüncü Moment Durumu):} $E[X_1^4] < \infty$ olsun.
\[
  E[(\bar{X}_n)^4] = \frac{1}{n^4} E\!\left[\left(\sum X_i\right)^4\right].
\]
$\sum X_i^4$ terim beklentisi $nE[X_1^4]$; çapraz terimler bağımsızlıkla yok olur;
$E[(X_i X_j)^2]$ terimleri $n(n-1)/2$ adet, her biri $E[X_1^2]^2$:
\[
  E[(\bar{X}_n)^4] = \frac{nE[X_1^4] + 3n(n-1)(E[X_1^2])^2}{n^4} \leq \frac{C}{n^2}.
\]
Markov ile $\varepsilon > 0$: $P(|\bar{X}_n| > \varepsilon) \leq C/(n^2\varepsilon^4)$.
$\sum_n P(|\bar{X}_n| > \varepsilon) \leq C/\varepsilon^4 \sum 1/n^2 < \infty$.
Borel-Cantelli 1: $P(|\bar{X}_n| > \varepsilon \text{ i.o.}) = 0$ her $\varepsilon > 0$ için.
Yani $\bar{X}_n \to 0$ n.k.

\textbf{Adım 2 (Genel Durum — Kolmogorov Kırpma):}
$E[|X_1|] < \infty$ için kırpılmış dizi $Y_n = X_n \mathbf{1}_{|X_n|\leq n}$ kullanılır.
Kolmogorov ölçütü: $\sum \mathrm{Var}(Y_n)/n^2 < \infty$ (Böl.~\ref{chp:olcum_teorisi}'deki
Kronecker lemması). Ayrıntılar Billingsley~\cite{Billingsley1995}, §22 veya Durrett~\cite{Durrett2019}, §2.4'te.
\end{proof}

\begin{teorem}[GBSK $\leftrightarrow$ ZBSK İlişkisi]
\label{thm:gbsk_zbsk}
GBSK, ZBSK'yi içerir: n.k.\ yakınsama $\Rightarrow$ olasılıkla yakınsama
(Teorem~\ref{thm:yaklinsama_hiyerarsi}). Tersi genel olarak yanlış.
\end{teorem}

\begin{ornek}[GBSK Uygulaması — Monte Carlo Hassasiyeti]
$g : [0,1] \to \mathbb{R}$ integrallenebilir. GBSK:
$\frac{1}{n}\sum_{i=1}^n g(U_i) \xrightarrow{\text{n.k.}} \int_0^1 g(x)\,dx$.
Bu, simülasyon çıktısının yalnızca "büyük ihtimalle" değil,
neredeyse tüm örnek yollarında kesinlikle doğru değere yakınsadığını garanti eder.
\end{ornek}

\begin{karsiornek}[GBSK'nın Beklenti Koşulunu Kaldıramazsınız]
$X_i \sim \mathrm{Cauchy}(0,1)$ bağımsız: $E[|X_1|] = \infty$.
$\bar{X}_n$, Cauchy(0,1) dağılımını korur — $0$'a yakınsamaz.
$E[|X_1|] < \infty$ koşulu vazgeçilmezdir.
\end{karsiornek}

\begin{mikroozet}
ZBSK: $\bar{X}_n \xrightarrow{P} \mu$ ($E[|X_1|] < \infty$, i.i.d.).
GBSK: $\bar{X}_n \xrightarrow{\text{n.k.}} \mu$ (aynı koşul, daha güçlü sonuç).
\end{mikroozet}

% =====================================================================
\section{Kolmogorov'un 0-1 Kanunu}
\label{sec:01_kanunu}
% =====================================================================

\begin{tanim}[Kuyruk $\sigma$-Cebiri]
\label{def:kuyruk_sigma}
$\{X_n\}$ rasgele değişkenler. $\mathcal{F}_n = \sigma(X_1, \ldots, X_n)$,
$\mathcal{T}_n = \sigma(X_{n+1}, X_{n+2}, \ldots)$. \textbf{Kuyruk $\sigma$-cebiri}:
\[
  \mathcal{T} = \bigcap_{n=1}^\infty \mathcal{T}_n.
\]
$\mathcal{T}$'nin olayları "sonlu sayıda $X_i$'yi değiştirmekten etkilenmeyen" olaylardır.
Örnekler: $\{\lim_n \bar{X}_n \text{ var}\}$, $\{\limsup_n X_n > c\}$, $\{\sum X_n \text{ yakınsar}\}$.
\end{tanim}

\begin{teorem}[Kolmogorov'un 0-1 Kanunu]
\label{thm:01_kanunu}
$X_1, X_2, \ldots$ bağımsız rasgele değişkenler ise kuyruk $\sigma$-cebirindeki
her $A \in \mathcal{T}$ olayı için $P(A) \in \{0, 1\}$.
\end{teorem}

\begin{proof}
\textbf{Adım 1:} Her $A \in \mathcal{T}$ ve her $n$ için $A \in \mathcal{T}_n = \sigma(X_{n+1}, X_{n+2}, \ldots)$,
dolayısıyla $A$, $\mathcal{F}_n = \sigma(X_1, \ldots, X_n)$'den bağımsızdır.

\textbf{Adım 2:} $\pi$-sistemi argümanı. $\mathcal{F}_\infty = \sigma(\bigcup_n \mathcal{F}_n)$.
Adım 1'deki bağımsızlık, $A$'nın $\mathcal{F}_\infty$'dan da bağımsız olduğunu gösterir
($\pi$-sistemi genişleme: $\bigcup_n \mathcal{F}_n$ bir $\pi$-sistemi; bağımsızlık üzerinde
$\lambda$-sistemi tutarlılığı ile Dynkin genişler).

\textbf{Adım 3:} $A \in \mathcal{T} \subseteq \mathcal{F}_\infty$ (her kuyruk olayı $\mathcal{F}_\infty$'dadır).
$A$, $\mathcal{F}_\infty$'dan bağımsız ve $A \in \mathcal{F}_\infty$:
$P(A) = P(A \cap A) = P(A) \cdot P(A)$.
Yani $P(A)^2 = P(A)$; çözüm: $P(A) \in \{0, 1\}$.
\end{proof}

\begin{ornek}[0-1 Kanunu Uygulamaları]
$\{X_n\}$ bağımsız rasgele değişkenler:
\begin{itemize}
  \item $P(\lim_n \bar{X}_n \text{ var}) \in \{0,1\}$. (GBSK bu olasılığın 1 olduğunu söyler.)
  \item $P(\sum X_n \text{ yakınsar}) \in \{0,1\}$.
  \item $P(\limsup_n X_n = +\infty) \in \{0,1\}$.
  \item $P(\limsup_n X_n/\sqrt{n\ln\ln n} = \sqrt{2}) \in \{0,1\}$. (Yinelenmiş logaritma kanunu.)
\end{itemize}
\end{ornek}

\begin{kritiksorgu}
0-1 Kanunu "olasılık ya 0 ya da 1" diyor. Bu deterministik mi?
Hayır — bu olaylar gerçek anlamda belirsiz olabilir; ancak birden fazla bağımsız
deneyin asimptotik davranışını yöneten olaylar, her örnek yolunu etkilese de
bütünleşik olarak keskin sınıflandırılabilir.
\end{kritiksorgu}

% =====================================================================
\section{Büyük Sapmalar: Giriş}
\label{sec:buyuk_sapmalar}
% =====================================================================

\begin{motivasyon}[Neden Büyük Sapmalar?]
ZBSK: $P(|\bar{X}_n - \mu| > \varepsilon) \to 0$. Chebyshev bize $O(1/n)$ hız verir.
Ama gerçek azalma çoğunlukla \emph{üstel}: $P(\bar{X}_n - \mu > \varepsilon) \approx e^{-nI(\varepsilon)}$.
Bu üstel azalma hızını veren $I(\varepsilon)$ fonksiyonu "oran fonksiyonu"dur.
\end{motivasyon}

\begin{teorem}[Chernoff Sınırı]
\label{thm:chernoff}
$X_1, \ldots, X_n$ i.i.d., $M(t) = E[e^{tX_1}]$ mevcut. Her $a > \mu = E[X_1]$ için
\[
  P(\bar{X}_n \geq a) \leq e^{-n \sup_{t \geq 0}(ta - \ln M(t))} = e^{-n \Lambda^*(a)},
\]
burada $\Lambda^*(a) = \sup_{t \geq 0}(ta - \ln M(t)) > 0$: \textbf{oran fonksiyonu}.
\end{teorem}

\begin{proof}
Her $t > 0$ için Markov ($e^{t\sum X_i} \geq e^{tna}$ olayından):
\[
  P(\bar{X}_n \geq a) = P\!\left(e^{t\sum X_i} \geq e^{tna}\right)
  \leq \frac{E[e^{t\sum X_i}]}{e^{tna}} = \frac{M(t)^n}{e^{tna}} = e^{n(\ln M(t) - ta)}.
\]
$t$'ye göre minimize etmek: $\inf_{t>0} e^{n(\ln M(t)-ta)} = e^{-n\sup_t(ta - \ln M(t))} = e^{-n\Lambda^*(a)}$.
\end{proof}

\begin{ornek}[Binom için Chernoff]
$X_i \sim \mathrm{Bernoulli}(p)$, $a > p$.
$\ln M(t) = \ln(q + pe^t)$.
$\Lambda^*(a) = a\ln(a/p) + (1-a)\ln((1-a)/q)$ (KL ıraksaması $\mathrm{Binom}$'dan $\mathrm{Binom}$'a).
$P(\bar{X}_n \geq a) \leq \exp(-n\,\mathrm{KL}(\mathrm{Ber}(a)\|\mathrm{Ber}(p)))$.
\end{ornek}

\begin{hocanotu}
Büyük sapmalar teorisi (Cramér teoremi, Varadhan'ın genel teorisi) bu kitabın
kapsamını aşar; ancak Chernoff sınırı istatistik ve algoritmik uygulamalarda
(öğrenme teorisi, randomize algoritmalar) sıkça kullanılır.
Derinlemesine okuma için Dembo-Zeitouni~\cite{Dembo1998}.
\end{hocanotu}

% ---- Disiplinlerarası Bağlantı (TYMM) ----
\begin{kopru}[Disiplinlerarası — Ekonomi, Fizik, Bilgisayar Bilimi]
\textbf{Ekonomi:} Büyük Sayılar Kanunu, ankete dayalı istatistiksel tahminlerin
(seçim anketi, tüketici fiyat endeksi) hangi koşulda güvenilir olduğunu açıklar:
gözlem sayısı $n$ büyüdükçe örneklem ortalaması nüfus ortalamasına yaklaşır.\\
\textbf{Fizik:} Ergodik teori (zamana göre ortalama = uzaya göre ortalama)
BSK'nın dinamik sistemler dilidir; termal dengedeki gaz molekülleri bu prensiple modellenir.\\
\textbf{Bilgisayar bilimi:} Randomize algoritmaların analizi BSK'ya dayanır:
$n$ tekrarlı bağımsız deneme, ortalama davranışı $O(1/\sqrt{n})$ hassasiyetle tahmin eder.
\end{kopru}

% ---- Öz-Yansıtma ----
\begin{ozyansima}
\begin{itemize}
  \item ZBSK ile GBSK arasındaki matematiksel fark nedir? Pratikte bu fark önemli mi?
  \item Borel-Cantelli lemmalarından birinin bağımsızlık koşulunu gerektirmesi,
        diğerinin gerektirmemesi hakkında ne düşünüyorsunuz?
  \item 0-1 Kanunu size "belirsizliğin çift değerli" olduğunu söylüyor;
        sezginizle bu nasıl bağdaşıyor?
  \item Chernoff sınırının Chebyshev'den daha güçlü olduğunu bir örnekle gösterin.
\end{itemize}
\end{ozyansima}

% ---- Köprü Notu ----
\begin{kopru}
\textbf{Önceki bölümler:} Bölüm~\ref{chp:beklenti} (Chebyshev, DCT) ve
Bölüm~\ref{chp:olcum_teorisi} (MYT) bu bölümün ispat araçlarını sağladı.\\
\textbf{Sonraki bölüm:} Bölüm~\ref{chp:donusumler} — Dönüşümler;
ardından Bölüm~\ref{chp:merkezi_limit} — BSK'nın öte boyutu:
$\bar{X}_n$ ne kadar hızlı yaklaşır? Cevap: $\sqrt{n}$ ölçeğinde Merkezi Limit Teoremi.\\
\textbf{İstatistik kısmı:} Bölüm~\ref{chp:istatistik_temelleri}'de büyük örneklem
tutarlılığı BSK çerçevesinde yeniden ele alınacak.
\defterbaglantisi{bolum07\_buyuk\_sayilar.ipynb}{%
  BSK simülasyonu, yakınsama hızı görselleştirmesi, Chernoff sınırı karşılaştırması}
\end{kopru}

% =====================================================================
\section{Yinelenmiş Logaritma Kanunu}
\label{sec:yil_kanunu}
% =====================================================================

\begin{motivasyon}[Büyük Sayılar Ötesinde]
GBSK: $\bar{X}_n \to \mu$ n.k. Ama ne \emph{kadar} hızlı yaklaşıyor?
$\sqrt{n}$ ölçeğinde Merkezi Limit Teoremi normal dağılım verir.
Daha ince soru: n.k.\ her örnek yolunun salınım büyüklüğü nedir?
Yinelenmiş Logaritma Kanunu (YLK) bu soruyu tam olarak yanıtlar.
\end{motivasyon}

\begin{teorem}[Kolmogorov'un Yinelenmiş Logaritma Kanunu]
\label{thm:ylk}
$X_1, X_2, \ldots$ i.i.d., $E[X_1] = 0$, $E[X_1^2] = \sigma^2 \in (0,\infty)$.
$S_n = X_1 + \cdots + X_n$ ise
\[
  P\!\left(\limsup_{n\to\infty} \frac{S_n}{\sqrt{2\sigma^2 n \ln\ln n}} = 1\right) = 1.
\]
Eşdeğer olarak: n.k.\ her örnek yolunda
$\limsup_{n\to\infty} S_n/\sqrt{2\sigma^2 n\ln\ln n} = 1$
ve $\liminf_{n\to\infty} S_n/\sqrt{2\sigma^2 n\ln\ln n} = -1$.
\end{teorem}

\begin{proof}[İspat Taslağı]
İki yönlü ispat gerekir.

\textbf{Üst sınır ($\limsup \leq 1$ n.k.):}
Sabit $b > 1$. Geometrik alt dizi $n_k = \lfloor b^k \rfloor$ seçin.
$A_k = \{S_{n_k} > (1+\varepsilon)\sqrt{2\sigma^2 n_k \ln\ln n_k}\}$.
Olağan sapma tahmini ile $P(A_k) \leq C/k^{1+\varepsilon}$.
$\sum P(A_k) < \infty$: Borel-Cantelli 1 ile $P(A_k \text{ i.o.}) = 0$.
Alt diziler arası maksimum kontrolü tüm $n$'ye genişler.

\textbf{Alt sınır ($\limsup \geq 1$ n.k.):}
Uygun $\varepsilon_k \to 0$ ile olaylar
$B_k = \{S_{n_{k+1}} - S_{n_k} > (1-\varepsilon)\sqrt{2\sigma^2(n_{k+1}-n_k)\ln\ln n_{k+1}}\}$
tanımlanır. $\{B_k\}$ asimptotik bağımsız; $\sum P(B_k) = \infty$.
Bağımsız Borel-Cantelli: $P(B_k \text{ i.o.}) = 1$.
Tam ispat: Hartman-Wintner (1941), bkz. Billingsley~\cite{Billingsley1995}, §12.
\end{proof}

\begin{ornek}[YLK'nın Yorumu]
$X_i \sim \Normal(0,1)$, $\sigma^2 = 1$. YLK:
\[
  \limsup_{n\to\infty} \frac{S_n}{\sqrt{2n\ln\ln n}} = 1 \quad \text{n.k.}
\]
$n = 10^6$: $\sqrt{2 \cdot 10^6 \cdot \ln\ln 10^6} = \sqrt{2\times10^6 \times 2.97} \approx 2439$.
$S_n$ sonsuz çok kez $\approx 2439$'a ulaşır, ama asla $2439 \cdot (1+\varepsilon)$'ı
sonsuz kez geçmez.
\end{ornek}

\begin{hataanalizi}
YLK ve MLT'yi karıştırmak yaygın hatadır.
MLT: $S_n/\sqrt{n\sigma^2} \xrightarrow{d} \Normal(0,1)$ — dağılımsal limit.
YLK: $\limsup S_n/\sqrt{2n\sigma^2\ln\ln n} = 1$ n.k. — yol bazında keskin sınır.
YLK, MLT'den \emph{daha güçlü} bir sonuçtur; her bir örnek yolu hakkında bilgi verir.
\end{hataanalizi}

\begin{mikroozet}
YLK: $S_n$'nin n.k.\ salınım büyüklüğü tam olarak $\sqrt{2\sigma^2 n \ln\ln n}$ ölçeğindedir.
GBSK ($S_n = o(n)$) ve MLT ($S_n = O(\sqrt{n})$) arasına düşer: $S_n = O(\sqrt{n\ln\ln n})$ n.k.
\end{mikroozet}

% =====================================================================
\section{Yenileme Teorisi Girişi}
\label{sec:yenileme_teorisi}
% =====================================================================

\begin{tanim}[Yenileme Süreci]
\label{def:yenileme_sureci}
$X_1, X_2, \ldots$ i.i.d., $X_i > 0$ n.k., $E[X_1] = \mu < \infty$.
$X_i$: arası zamanlar. $S_0 = 0$, $S_n = X_1 + \cdots + X_n$.
\textbf{Yenileme sayacı}: $N(t) = \max\{n \geq 0 : S_n \leq t\}$.
$\{N(t), t \geq 0\}$: \textbf{yenileme süreci} (renewal process).
\end{tanim}

\begin{teorem}[Temel Yenileme Teoremi]
\label{thm:temel_yenileme}
$E[X_1] = \mu \in (0,\infty)$ ise
\[
  \frac{N(t)}{t} \xrightarrow{\text{n.k.}} \frac{1}{\mu} \quad (t \to \infty).
\]
Daha kesin: $E[N(t)]/t \to 1/\mu$ ve $\mathrm{Var}(N(t))/t \to \sigma^2/\mu^3$
($\sigma^2 = \mathrm{Var}(X_1)$).
\end{teorem}

\begin{proof}
$N(t) \geq n \Leftrightarrow S_n \leq t$. Dolayısıyla $S_{N(t)} \leq t < S_{N(t)+1}$.
Her iki tarafı $N(t)$'ye böl:
\[
  \frac{S_{N(t)}}{N(t)} \leq \frac{t}{N(t)} < \frac{S_{N(t)+1}}{N(t)}.
\]
GBSK: $S_{N(t)}/N(t) \xrightarrow{\text{n.k.}} \mu$ ($N(t) \to \infty$ çünkü $\mu < \infty$).
Sıkıştırma: $t/N(t) \xrightarrow{\text{n.k.}} \mu$, yani $N(t)/t \xrightarrow{\text{n.k.}} 1/\mu$.
\end{proof}

\begin{ornek}[Makine Bakım Modeli]
Bir makine $X_i \sim \Exp(\lambda)$ ömrü sonunda yenileniyor.
$E[X_i] = 1/\lambda$. Temel Yenileme: $N(t)/t \to \lambda$ n.k.
$t$ birim zamanda ortalama $\lambda t$ yenileme yapılır.
\end{ornek}

\begin{teorem}[Blackwell'in Yenileme Teoremi]
\label{thm:blackwell}
$X_i$ aritmetik-olmayan (non-arithmetic) dağılımlı ise: her $h > 0$ için
\[
  E[N(t+h) - N(t)] \to \frac{h}{\mu} \quad (t \to \infty).
\]
Yani uzun vadede $[t, t+h]$ aralığındaki beklenen yenileme sayısı $h/\mu$'ya yaklaşır:
sistem "kararlı duruma" (steady state) girer.
\end{teorem}

\begin{disiplinlerarasibaglan}
\textbf{Güvenilirlik Mühendisliği:} Yenileme süreci, ekipman ömürleri ve
bakım planlamasının temel modelidir: $N(t)$ arızaların kümülatif sayısıdır.\\
\textbf{Kuyruk Teorisi:} M/G/1 kuyruğunda servis süreleri yenileme sürecini oluşturur;
Temel Yenileme Teoremi ortalama bekleme süresi hesaplamasında kullanılır.\\
\textbf{Ekonomi:} Stok yönetiminde sipariş dönemleri yenileme yapısı taşır.
\end{disiplinlerarasibaglan}

% =====================================================================
\section{Durağan Diziler için Ergodik Teorem}
\label{sec:ergodik_teorem}
% =====================================================================

\begin{tanim}[Durağan Dizi]
\label{def:duragan_dizi}
$\{X_n\}$ rasgele değişkenler dizisi \textbf{(geniş anlamda) durağan}
(stationary) denir, eğer her $k, n$ ve $h$ için
$(X_k, X_{k+1}, \ldots, X_{k+n})$ ile $(X_{k+h}, X_{k+h+1}, \ldots, X_{k+h+n})$
aynı dağılıma sahipse.

Özel durumlar: i.i.d.\ diziler, durağan Markov zincirleri.
\end{tanim}

\begin{tanim}[Kaydırma Dönüşümü ve Ergodiklik]
\label{def:ergodiklik}
$\Omega$'da \textbf{kaydırma dönüşümü} $\theta: \theta(\omega)_n = \omega_{n+1}$.
Dizi $\theta$-uyumlu: $X_n(\omega) = X_0(\theta^n \omega)$.
$A \in \mathcal{F}$ \textbf{değişmez} (invariant) olay: $\theta^{-1}A = A$.
$P$ ölçümü \textbf{ergodik}: her değişmez olay $A$ için $P(A) \in \{0,1\}$.
\end{tanim}

\begin{teorem}[Birkhoff'un Ergodik Teoremi]
\label{thm:birkhoff}
$\{X_n\}$ durağan ve ergodik (ölçüm $P$ ergodik), $E[|X_0|] < \infty$. O zaman
\[
  \frac{1}{n}\sum_{k=0}^{n-1} X_k \xrightarrow{\text{n.k.}} E[X_0].
\]
\end{teorem}

\begin{proof}[İspat Notu]
Bağımsız diziler için (i.i.d.) bu Kolmogorov GBSK'dır.
Genel ergodik durum için ispat daha derindir: Maximal Ergodik Lemma (Maximal Lemma)
ve Garsia'nın kombinatoryal argümanı kullanılır.
Referans: Durrett~\cite{Durrett2019}, §7.2; Billingsley~\cite{Billingsley1995}, §24.
\end{proof}

\begin{ornek}[Ergodik Markov Zinciri]
$\{Z_n\}$ sonlu durum uzaylı, geri dönüşümlü ve pozitif tekrarlı Markov zinciri;
$\pi$ durağan dağılımı. $f: S \to \mathbb{R}$ her duruma bir değer atar.
Birkhoff: $\frac{1}{n}\sum_{k=0}^{n-1} f(Z_k) \xrightarrow{\text{n.k.}} \sum_{s\in S} f(s)\pi(s) = E_\pi[f]$.

Yani "zamansal ortalama = uzamsal ortalama" — klasik ergodiklik eşitliği.
\end{ornek}

\begin{kritiksorgu}
Ergodik teorem i.i.d.\ koşulunu gevşetiyor. Hangi koşulun gerektiğini tartışın:
bağımsızlık mı, durağanlık mı, yoksa ergodiklik mi? Ergodikliğin 0-1 Kanunu
ile ilişkisini açıklayın.
\end{kritiksorgu}

% =====================================================================
\section{Yavaş Değişen Fonksiyonlar ve Düzenli Değişim}
\label{sec:duzenli_degisim}
% =====================================================================

\begin{tanim}[Yavaş değişen fonksiyon]\label{def:yavas_degisen}
$L:(0,\infty)\to(0,\infty)$ fonksiyonu \emph{yavaş değişen} (\emph{slowly varying}) denir, her $c>0$ için
\[
  \lim_{x\to\infty}\frac{L(cx)}{L(x)}=1.
\]
\end{tanim}

\begin{ornek}[Yavaş değişen fonksiyon örnekleri]\label{ex:yavas_ornekler}
\begin{enumerate}
  \item Sabitler: $L(x)=c>0$.
  \item $L(x)=\log x$, $L(x)=(\log\log x)^2$, $L(x)=\log(1+x^2)$.
  \item İterated logarithm: $L(x)=\log_k(x)$ ($k$'ıncı yinelenmiş logaritma).
\end{enumerate}
\emph{Düzenli değişen}: $R(x)=x^\alpha L(x)$ ($L$ yavaş değişen, $\alpha\in\mathbb{R}$ indeks).
\end{ornek}

\begin{teorem}[Karamata temsil teoremi]\label{thm:karamata}
$L$ yavaş değişen $\iff$ $L(x)=c(x)\exp\!\left(\int_1^x \frac{\varepsilon(t)}{t}dt\right)$,
burada $c(x)\to c>0$ ve $\varepsilon(t)\to0$.
\end{teorem}

\begin{teorem}[Potter sınırları]\label{thm:potter}
$L$ yavaş değişen. Her $\delta>0$ için $x_0>0$ mevcut:
$x,y\geq x_0$ ise
\[
  \frac{L(y)}{L(x)} \leq 2\max\!\left\{\left(\frac{y}{x}\right)^\delta, \left(\frac{x}{y}\right)^\delta\right\}.
\]
\end{teorem}

\begin{ornek}[Ağır kuyruklu dağılımlar]\label{ex:agir_kuyruk}
$P(X>x)\sim x^{-\alpha}L(x)$ ($\alpha>0$, $L$ yavaş değişen): \emph{düzenli değişen kuyruk}.
\begin{itemize}
  \item $\alpha>2$: $E[X^2]<\infty$ — standart MLT geçerli.
  \item $1<\alpha\leq2$: $E[X]<\infty$ ama $E[X^2]=\infty$ — kararlı dağılımlara yakınsama, ölçek $n^{1/\alpha}$.
  \item $0<\alpha\leq1$: $E[X]=\infty$ — büyük sayılar kanunu çöker.
\end{itemize}
Pareto$(\alpha)$, Cauchy $(\alpha=1)$, Student-$t_\nu$ ($\nu=\alpha$) bu sınıftadır.
\end{ornek}

\begin{hataanalizi}
"Büyük sayılar kanunu her zaman geçerlidir" yanlıştır. $E[|X|]=\infty$ durumunda ZBSK ve GBSK çöker: Cauchy dağılımında örneklem ortalaması hiçbir değere yakınsamaz, her örneklem için farklı bir değer verir. Bu, finansta "ağır kuyruklu" kayıpların klasik istatistikle modellenememesinin matematiksel gerekçesidir.
\end{hataanalizi}

% =====================================================================
\section{Büyük Sapmalar İçin Üst Sınırlar}
\label{sec:buyuk_sapma_sinir}
% =====================================================================

\begin{teorem}[Markov eşitsizliği — hatırlatma]\label{thm:markov_hatirlatma}
$X\geq0$, $a>0$: $P(X\geq a)\leq E[X]/a$.
\end{teorem}

\begin{teorem}[Chebyshev eşitsizliği]\label{thm:chebyshev_hatirlatma}
$E[X]=\mu$, $\Var(X)=\sigma^2<\infty$, $k>0$:
$P(|X-\mu|\geq k\sigma)\leq 1/k^2$.
\end{teorem}

\begin{teorem}[Chernoff sınırı]\label{thm:chernoff}
$X_1,\ldots,X_n\iid$, MGF $M(\lambda)=E[e^{\lambda X}]$ $\lambda_0>0$ için sonlu. $S_n=\sum_{i=1}^n X_i$. Her $t>E[X]$ için:
\[
  P(S_n\geq nt) \leq e^{-n\sup_{\lambda\geq0}[\lambda t - \log M(\lambda)]}.
\]
Üstel $[\lambda t - \log M(\lambda)]$ büyük sapma hızıdır (\emph{rate function} $I(t)$).
\end{teorem}

\begin{proof}
$P(S_n\geq nt)=P(e^{\lambda S_n}\geq e^{n\lambda t})\leq e^{-n\lambda t}E[e^{\lambda S_n}]=e^{-n\lambda t}[M(\lambda)]^n=e^{-n[\lambda t-\log M(\lambda)]}$.
$\lambda\geq0$'da supremum alınır.
\end{proof}

\begin{ornek}[Hoeffding eşitsizliği]\label{ex:hoeffding}
$X_i\in[a_i,b_i]$ sınırlı bağımsız, $E[X_i]=0$. $P(S_n\geq t)\leq\exp\!\left(-\frac{2t^2}{\sum(b_i-a_i)^2}\right)$.
Özellikle $a_i=0$, $b_i=1$, $t=n\varepsilon$: $P(\bar X_n\geq\varepsilon)\leq e^{-2n\varepsilon^2}$.
\end{ornek}

\begin{disiplinlerarasibaglan}
\textbf{Makine Öğrenmesi — PAC öğrenimi:} Hoeffding eşitsizliği, bir öğrenme algoritmasının eğitim hatasının test hatasına $\varepsilon$ içinde yakınsadığını yüksek olasılıkla garanti eden PAC (\emph{probably approximately correct}) öğrenimi sınırının temelidir.

\textbf{Finansal Risk:} Düzenli değişen kuyruklar (Pareto, $t_\nu$) piyasa çöküşlerini modellemek için kullanılır; Basle III sermaye yeterliliği hesaplamalarında VaR (değer-risk) bu dağılımlara dayanır.
\end{disiplinlerarasibaglan}

\begin{mikroozet}
Durağan + ergodik dizilerde $\frac{1}{n}\sum_{k=0}^{n-1}X_k \to E[X_0]$ n.k.
Bu, GBSK'nın i.i.d.\ ötesine genellemesidir; Markov zincirleri ve zaman serisi
analizinin matematiksel temelidir.
Düzenli değişen kuyruklar ($P(X>x)\sim x^{-\alpha}L(x)$): $\alpha\leq1$'de BSK çöker.
Chernoff/Hoeffding: üstel büyük sapma sınırları; makine öğrenmesinde PAC sınırlarının temeli.
\end{mikroozet}

% =====================================================================
%  ALIŞTIRMALAR
% =====================================================================

\cozumlualistirmalar

\begin{alistirma}[Al.7.1]{A}{Yakınsama Türlerini Tanımak}
$X_n = n^{-1/2}Z$ ($Z \sim \Normal(0,1)$, sabit). Gösterin:
(a) $X_n \xrightarrow{\text{n.k.}} 0$. (b) $X_n \xrightarrow{P} 0$. (c) $X_n \xrightarrow{d} 0$.
Bu örnek hangi tipik yanılgıyı düzeltiyor?
\end{alistirma}
\alcozum{%
$X_n(\omega) = Z(\omega)/\sqrt{n} \to 0$ her $\omega$'da (sabit $Z$, deterministik yakınsama).
(a)--(c) üçü de trivyal olarak sağlanır.\\
Düzeltilen yanılgı: Burada $X_n$ "aynı $Z$'yi ölçekliyor"; dizi gerçekten bağımsız değil.
Bağımsız $X_n$'lerin neredeyse kesin yakınsaması için Borel-Cantelli gerekir.}

\begin{alistirma}[Al.7.2]{A}{Borel-Cantelli Uygulaması}
$X_1, X_2, \ldots$ bağımsız, $P(X_n > n) = 1/n^2$.
$P(X_n > n \text{ sonsuz çok } n \text{ için}) = ?$
\end{alistirma}
\alcozum{%
$\sum_{n=1}^\infty P(X_n > n) = \sum 1/n^2 = \pi^2/6 < \infty$.
Borel-Cantelli 1: $P(X_n > n \text{ i.o.}) = 0$.
Yani yalnızca sonlu sayıda $n$ için $X_n > n$ gerçekleşir.}

\begin{alistirma}[Al.7.3]{B}{ZBSK ve GBSK Karşılaştırması}
$X_1, X_2, \ldots$ i.i.d.\ $\mathrm{Exp}(1)$. $\bar{X}_n$'nin hangi değere yakınsadığını
belirtin. (a) Chebyshev ZBSK ile $P(|\bar{X}_n - 1| > \varepsilon)$ için $n$ sayısına bağlı
üst sınır verin. (b) Bu yakınsamanın GBSK anlamında da geçerli olduğunu söyleyin.
\end{alistirma}
\alcozum{%
$E[X_1] = 1$, $\mathrm{Var}(X_1) = 1$.\\
(a) Chebyshev: $P(|\bar{X}_n - 1| > \varepsilon) \leq 1/(n\varepsilon^2)$.
$\varepsilon = 0.1$ için $n = 100$ yeterli: $P(\ldots) \leq 1$.
(İyi bir üst sınır değil ama metodoloji doğru; Chernoff çok daha iyi.)\\
(b) $E[|X_1|] = 1 < \infty$, i.i.d.; GBSK: $\bar{X}_n \xrightarrow{\text{n.k.}} 1$.}

\begin{alistirma}[Al.7.4]{B}{Chernoff Sınırı}
$X_i \sim \mathrm{Bernoulli}(0.5)$ bağımsız, $n = 100$, $a = 0.6$.
(a) Chebyshev ile $P(\bar{X}_{100} \geq 0.6)$ için üst sınır verin.
(b) Chernoff ile üst sınır verin ve karşılaştırın.
\end{alistirma}
\alcozum{%
(a) Chebyshev: $P(\bar{X}_n \geq 0.6) \leq P(|\bar{X}_n-0.5|\geq 0.1) \leq \mathrm{Var}(\bar{X}_n)/(0.1)^2
= (0.25/100)/0.01 = 0.25$.\\
(b) Chernoff: $\Lambda^*(0.6) = 0.6\ln(0.6/0.5)+0.4\ln(0.4/0.5) = 0.6\ln(1.2)+0.4\ln(0.8)$.
$\approx 0.6(0.1823)+0.4(-0.2231) = 0.1094-0.0892 = 0.0202$.
$P(\bar{X}_{100}\geq 0.6)\leq e^{-100\times0.0202} = e^{-2.02} \approx 0.133$.
Chebyshev: 0.25; Chernoff: 0.133. Gerçek değer $\approx 0.028$ (Chernoff daha yakın).}

\begin{alistirma}[Al.7.4b]{B}{Bağımsızlık Olmadan BSK — Ergodik Teorem}
$(X_n)_{n\geq1}$ \emph{durağan ergodik} dizi; $E[|X_1|]<\infty$, $\mu=E[X_1]$.
(a) Birkhoff ergodik teoremi ne söyler? Güçlü BSK ile nasıl kıyaslanır?
(b) Markov zinciri $X_n$ durağan dağılım $\pi$ ile ergodikse $\frac{1}{n}\sum_{k=1}^n f(X_k)\to\sum_x f(x)\pi(x)$ n.k. Bunu simülasyon bağlamında yorumlayın.
(c) $X_1,X_2,\ldots$ i.i.d.\ ama simetrik Cauchy dağılımlı. $\bar X_n$ ne yapar? Neden BSK çalışmaz?
\end{alistirma}
\alcozum{%
(a) \textbf{Birkhoff ergodik teoremi}: Durağan ergodik $(X_n)$ için $\bar X_n=\frac{1}{n}\sum_{k=1}^n X_k\xrightarrow{\text{n.k.}}\mu=E[X_1]$. Güçlü BSK, i.i.d.\ varsayımı altında ergodikliğin özel halidir; Birkhoff bağımlı ama durağan-ergodik dizilere geneller.

(b) Markov Monte Carlo (MCMC) temeli: Zinciri durağan dağılımdan başlatın (veya yeterince uzun ısın sonra) ve $\frac{1}{n}\sum f(X_k)$ ile $E_\pi[f]$'yi tahmin edin. Ergodiklik garantisi yakınsamayı sağlar; bağımsız örnekleme yerine geçer.

(c) Cauchy için $E[|X_1|]=\infty$ (beklenti var olmaz). Güçlü BSK'nın koşulu $E[|X_1|]<\infty$ ihlal edilir. Gerçekte $\bar X_n$ Cauchy($0,1$) dağılımlı kalır: $n$ büyüse de dağılım daralmaz. Büyük sayılar kanunları geçerli değil.}

\begin{alistirma}[Al.7.4c]{C}{Yavaş Değişen Fonksiyonlar ve Kuyruğun Ağırlığı}
$L$ yavaş değişen (slowly varying) demek: her $c>0$ için $L(cx)/L(x)\to1$ ($x\to\infty$).
(a) $L(x)=\log x$ ve $L(x)=(\log x)^2$'nin yavaş değişen olduğunu doğrulayın.
(b) $P(X>x)=L(x)/x^\alpha$ ($\alpha>0$) ise $X$'in ortalaması ne zaman sonlu?
(c) $\alpha=1$, $L(x)=1$ durumunda (Cauchy benzeri), $S_n/n\to?$ Limitin ne olduğunu belirtin.
\end{alistirma}
\alcozum{%
(a) $L(x)=\log x$: $\frac{\log(cx)}{\log x}=\frac{\log c+\log x}{\log x}=1+\frac{\log c}{\log x}\to1$. $\checkmark$ Benzer şekilde $(\log x)^2$ için de aynı argüman.

(b) $E[X]=\int_0^\infty P(X>x)dx=\int_0^1 dx+\int_1^\infty \frac{L(x)}{x^\alpha}dx$. İkinci integral $\int_1^\infty x^{-\alpha}dx$ ile karşılaştırılır (Potter lemması: $L(x)\leq Cx^\varepsilon$ yeterince büyük $x$'te). Sonlu olması için $\alpha>1$ gerekir ($\alpha=1$ sınır: $\log$ düzeltmesiyle ince analiz).

(c) $\alpha=1$, $P(X>x)\sim1/x$ ($L=1$): Cauchy dağılımı modeli. Stabil dağılım teorisi: $S_n/(n\log n)^{1/\alpha}$ için uygun normalleştirme gerekir; $S_n/n$ hiçbir sonlu limite yakınsamaz ($\bar X_n$ Cauchy kalır). BSK ve CLT her ikisi de çöker.}

% ---

\alistirmalar

\begin{alistirma}[Al.7.5]{A}{Yakınsama Türleri}
Her biri için örnek verin ya da imkânsız olduğunu gösterin:
(a) $X_n \xrightarrow{P} X$ ama $X_n \not\xrightarrow{\text{n.k.}} X$.
(b) $X_n \xrightarrow{d} c$ (sabit) ama $X_n \not\xrightarrow{P} c$.
\end{alistirma}
\alcozum{%
(a) Gezgin blok (Karşı Örnek~\ref{ex:yaklinsama_karsiornek}).\\
(b) İmkânsız: $X_n \xrightarrow{d} c$ (dejenere) $\Leftrightarrow$ $X_n \xrightarrow{P} c$.
Çünkü $P(|X_n-c|>\varepsilon) = P(X_n \leq c-\varepsilon) + P(X_n > c+\varepsilon) \to 0+0=0$.}

\begin{alistirma}[Al.7.6]{A}{0-1 Kanunu Uygulaması}
$X_1, X_2, \ldots$ bağımsız. Aşağıdaki olayların olasılığını belirtin:
(a) $A = \{\bar{X}_n \to 0\}$.
(b) $B = \{\sum_{n=1}^\infty X_n/n \text{ yakınsar}\}$.
(c) $C = \{\limsup_n X_n < \infty\}$.
\end{alistirma}
\alcozum{%
Hepsi kuyruk $\sigma$-cebirinde. 0-1 Kanunu: $P(A), P(B), P(C) \in \{0,1\}$.
Belirli dağılıma bağlı olmadan bu sonuç kesin — spesifik değer ayrıca hesaplanmalı.}

\begin{alistirma}[Al.7.7]{B}{Borel-Cantelli — Kuvvetli Kanun}
$X_1, X_2, \ldots$ bağımsız, $E[X_i] = \mu$, $\mathrm{Var}(X_i) = \sigma^2 < \infty$.
Borel-Cantelli 1'i kullanarak $\bar{X}_{n^2} \xrightarrow{\text{n.k.}} \mu$ olduğunu ispatlayın.
($n^2$'nin kare indeksler olduğuna dikkat: $\bar{X}$ dizisinin alt dizisi.)
\end{alistirma}
\alcozum{%
$\varepsilon > 0$. Chebyshev: $P(|\bar{X}_{n^2}-\mu|>\varepsilon) \leq \sigma^2/(n^2\varepsilon^2)$.
$\sum_n 1/n^2 < \infty$: $\sum_n P(|\bar{X}_{n^2}-\mu|>\varepsilon) < \infty$.
Borel-Cantelli 1: $P(|\bar{X}_{n^2}-\mu|>\varepsilon \text{ i.o.}) = 0$.
Her $\varepsilon>0$ için geçerli: $\bar{X}_{n^2} \xrightarrow{\text{n.k.}} \mu$.}

\begin{alistirma}[Al.7.8]{B}{Slutsky Uygulaması}
$X_n \xrightarrow{d} \Normal(0,1)$ ve $S_n^2 \xrightarrow{P} \sigma^2 > 0$ (örnek varyans).
$T_n = X_n / S_n$'nin limit dağılımını bulun.
\end{alistirma}
\alcozum{%
$Y_n = S_n \xrightarrow{P} \sigma$; $g(y) = 1/y$ sürekli ($y\neq0$).
Sürekli haritalama: $1/S_n \xrightarrow{P} 1/\sigma$.
Slutsky çarpım: $T_n = X_n \cdot (1/S_n) \xrightarrow{d} \Normal(0,1) \cdot (1/\sigma) = \Normal(0,1/\sigma^2)$.}

\begin{alistirma}[Al.7.9]{C}{GBSK — Kolmogorov Yeterli Koşulu}
$\{X_n\}$ bağımsız (mutlaka i.i.d.\ değil), $E[X_n] = 0$,
$\sum_{n=1}^\infty \mathrm{Var}(X_n)/n^2 < \infty$.
Kronecker lemmasını kullanarak $\bar{X}_n \xrightarrow{\text{n.k.}} 0$ olduğunu gösterin.
(\textit{İpucu:} $S_n = \sum_{k=1}^n X_k/k$'nin n.k.\ yakınsadığını Borel-Cantelli ile kurun;
ardından Kronecker: $S_n \to S$ ve $S_n = \sum X_k/k$ ise $\frac{1}{n}\sum_{k=1}^n X_k \to 0$.)
\end{alistirma}
\alcozum{%
\begin{ipucu}{1}
$Y_k = X_k/k$ tanımlayın; $\sum E[Y_k^2] = \sum \mathrm{Var}(X_k)/k^2 < \infty$.
$L^2$ yakınsaması + martingal argümanı veya Borel-Cantelli ile $\sum Y_k$ n.k.\ yakınsar.
\end{ipucu}
\textbf{Kronecker Lemması:} $\sum_{k=1}^\infty a_k/k$ yakınsarsa $\frac{1}{n}\sum_{k=1}^n a_k \to 0$.
$a_k = X_k$ uygulandığında: $\sum X_k/k$ yakınsıyor (yukarıdaki ipucu) $\Rightarrow$ $\bar{X}_n \to 0$ n.k.}

\begin{alistirma}[Al.7.10]{B}{Yinelenmiş Logaritma Kanunu — Uygulama}
$X_i \sim \Normal(0,1)$ i.i.d., $S_n = \sum_{i=1}^n X_i$.
(a) YLK'ya göre $\limsup_{n\to\infty} S_n/\sqrt{2n\ln\ln n}$'nin değerini verin.
(b) $n = 10^4$ için $\sqrt{2n\ln\ln n}$'yi sayısal olarak hesaplayın.
(c) "GBSK $\bar{X}_n \to 0$" ile "YLK $S_n/\sqrt{2n\ln\ln n} \to 1$" ifadelerinin
    çelişmediğini açıklayın.
\end{alistirma}
\alcozum{%
(a) YLK (Teorem~\ref{thm:ylk}): $\sigma^2 = 1$ olduğundan $\limsup = 1$ n.k.\\
(b) $n = 10^4$: $\ln n = \ln 10^4 = 4\ln 10 \approx 9.21$; $\ln\ln n = \ln 9.21 \approx 2.22$.
$\sqrt{2 \times 10^4 \times 2.22} \approx \sqrt{44400} \approx 210.7$.
Yani $S_{10^4}$ sonsuz kez $\approx 210.7$'ye ulaşır.\\
(c) $\bar{X}_n = S_n/n \to 0$ (GBSK); $S_n/\sqrt{2n\ln\ln n} \to 1$ (YLK).
$n^{-1} = o(n^{-1/2}(\ln\ln n)^{-1/2})$ olduğundan iki ifade farklı ölçek dilleridir.
$\bar{X}_n \to 0$ yavaş yavaş gerçekleşir; YLK tam ne kadar yavaş olduğunu söyler.}

\begin{alistirma}[Al.7.11]{B}{Yenileme Teorisi}
Bir makine parçası $X_i \sim \Exp(2)$ dağılımına sahip ömürler geçiriyor ($i = 1, 2, \ldots$).
Parça bozulduğunda hemen yenisiyle değiştiriliyor. $N(t)$ yenileme sayısı.
(a) Temel Yenileme Teoremi ile $\lim_{t\to\infty} N(t)/t$'yi bulun.
(b) Blackwell teoremini kullanarak $[100, 101]$ aralığındaki beklenen yenileme sayısını yaklaşık hesaplayın.
(c) $N(t)$'nin standart sapmasının $\sim \sqrt{\sigma^2 t/\mu^3}$ ölçeğinde olduğunu yorumlayın
    ($\sigma^2 = \mathrm{Var}(X_1) = 1/4$, $\mu = 1/2$).
\end{alistirma}
\alcozum{%
$\mu = E[X_1] = 1/2$, $\sigma^2 = \mathrm{Var}(X_1) = 1/4$ ($\Exp(2)$ için).\\
(a) $N(t)/t \xrightarrow{\text{n.k.}} 1/\mu = 2$. Uzun vadede dakikada 2 yenileme.\\
(b) Blackwell: $E[N(101)-N(100)] \approx h/\mu = 1/(1/2) = 2$.
$[100,101]$ aralığında beklenen 2 yenileme.\\
(c) $\sigma^2 t/\mu^3 = (1/4) \cdot t/(1/8) = 2t$.
$\mathrm{SD}(N(t)) \approx \sqrt{2t}$: büyük $t$'de $N(t)$ etrafında $\sqrt{2t}$ ölçeğinde salınır.}

\begin{alistirma}[Al.7.12]{C}{Ergodik Markov Zinciri}
$Z_n$ iki durumlı Markov zinciri, $S = \{0, 1\}$, geçiş matrisi
$P = \begin{pmatrix}1-\alpha & \alpha \\ \beta & 1-\beta\end{pmatrix}$
($0 < \alpha, \beta < 1$).
(a) Durağan dağılım $\pi$'yi bulun.
(b) $f(0) = 0$, $f(1) = 1$ ise Birkhoff Ergodik Teoremi'ne göre
    $\frac{1}{n}\sum_{k=0}^{n-1} f(Z_k)$'nın sınırını verin.
(c) Bu sonucu BSK ile karşılaştırın: bağımsızlık gerekiyor mu?
\end{alistirma}
\alcozum{%
(a) $\pi P = \pi$: $\pi_0(1-\alpha) + \pi_1 \beta = \pi_0$ ve $\pi_0 + \pi_1 = 1$.
$\pi_0 \alpha = \pi_1 \beta \Rightarrow \pi_1/\pi_0 = \alpha/\beta$.
$\pi_0 = \beta/(\alpha+\beta)$, $\pi_1 = \alpha/(\alpha+\beta)$.\\
(b) Zincir geri dönüşümlü ve pozitif tekrarlı (sonlu durum, $\alpha,\beta > 0$) $\Rightarrow$ ergodik.
Birkhoff: $\frac{1}{n}\sum f(Z_k) \xrightarrow{\text{n.k.}} E_\pi[f] = \pi_1 = \alpha/(\alpha+\beta)$.\\
(c) BSK: bağımsızlık + i.i.d.\ gerekir. Ergodik teorem: yalnızca durağanlık + ergodiklik.
Markov zinciri bağımlı olduğu hâlde "zamansal ortalama = uzamsal ortalama" geçerlidir.
Bağımsızlık, ergodikliğin özel bir durumudur.}

% =====================================================================
\endinput
